%% Supporting Information for the Water Resources Research submission.
%% Rewritten 2026-10-20 to match the aligned-policy evidence chain
%% (results/wrr_submission_ready/major_revision_20261020/). Built by
%% scripts/223_wrr_build.py alongside the article.

\documentclass[draft]{agujournal2025}

\usepackage{booktabs}
\usepackage{tabularx}
\newcolumntype{R}[1]{>{\raggedleft\arraybackslash}p{#1}}
\usepackage{microtype}
\ifdefined\XeTeXversion\else
  \usepackage{seqsplit}
  \let\wrrtexttt\texttt
  \renewcommand{\texttt}[1]{\wrrtexttt{\seqsplit{#1}}}
\fi
\setlength{\emergencystretch}{1.5em}
\usepackage{amsmath}

\renewcommand{\thefigure}{S\arabic{figure}}
\renewcommand{\thetable}{S\arabic{table}}
\renewcommand{\thesection}{S\arabic{section}}

%% Table bodies inherit the article's double spacing; switch it off as in the
%% main text so a wrapped cell does not become a staircase.
\newcommand{\tabletext}{%
  \footnotesize
  \setlength{\baselineskip}{10pt}%
  \setlength{\parindent}{0pt}%
}

\begin{document}

\journalname{Water Resources Research}

\title{Supporting Information for ``Hiding the peak, not missing the days:
support-based screening of incomplete river-temperature archives''}

\authors{AUTHOR ONE\affil{1}}
\affiliation{1}{AFFILIATION PENDING AUTHOR CONFIRMATION}
\authoraddr{Corresponding author details pending author confirmation.}
\authorrunninghead{Ignored.}
\titlerunninghead{Ignored.}

\maketitle

\section*{Contents of this file}

\begin{enumerate}
  \item Text S1 to S12
  \item Figure S1
  \item Tables S1 to S24
\end{enumerate}

\section*{Introduction}

This supporting information records what belongs to the record rather than to the
argument: data provenance, protocol history, the predictions that were written
down and not met, the runs that were discarded and why, and the reproduction path.
Nothing here is required to follow the article.

\section{Data sources}

Daily water-temperature series were assembled from nine public providers
(Table~\ref{tab:si-providers}); the external panel adds France's Hub'Eau and the
UK Environment Agency. Each provider's product, documented statistic, quality
codes, licence and acquisition route were re-verified against the provider's own
documentation on 2026-10-12
(\texttt{results/wrr\_submission\_ready/measurement\_provenance\_dictionary.csv}).
Eight of the eleven publish under a named open licence and none required a
data-use agreement, but LUBW reserves copyright on its raw values: those seven
networks are not redistributed here and are read locally through LUBW's own
service.

\begin{table}
\caption{Providers and their role in the study.}
\label{tab:si-providers}
\centering
\begin{tabular}{lll}
\hline
provider & country & role \\
\hline
ARSO & Slovenia & development \\
CHMI & Czechia & development \\
FOEN & Switzerland & development \\
GKD Bayern & Germany (Bavaria) & development \\
LUBW & Germany (Baden-W\"urttemberg) & development \\
NVE & Norway & development \\
Rijkswaterstaat & Netherlands & development \\
SYKE & Finland & development \\
USGS & United States & development and external \\
Hub'Eau & France & external \\
UK Environment Agency & United Kingdom & external \\
\hline
\end{tabular}
\end{table}

Station coordinates were assembled for all 457 stations: from the providers'
own catalogues already held in the repository for FOEN, CHMI, SYKE, LUBW, GKD
Bayern, NVE and Rijkswaterstaat and for nine USGS gauges; from the USGS NWIS
site service for the remaining United States gauges; and from ARSO's public
station feed. Every row of the coordinate table in the repository records its
source.

Daily 2\,m air temperature is NASA POWER (MERRA-2), 1 January 1981 to 31 August
2026, one point per station at the product's native resolution. Basin
characteristics for the United States subset are GAGES-II, joined on the
eight-digit gauge identifier.

\section{Quality control and the definition of a network}

A \emph{network} is a provider's group of stations on a common river system,
which is the unit at which donor stations are available to one another. Networks
are the clustering unit for every interval in the article, because station-years
within a network share weather, regulation and instrumentation. Two labels that
share a station are merged into one \emph{independence group} before any fold is
drawn, which is why 49 external labels form 41 validation groups.

Waterbody class was assigned per network from the provider's own naming and an
explicit map; seas, lakes, canals, estuaries and deltas are listed and excluded
rather than averaged into a river result. Reference series are never imputed: a
calendar year with any missing day contributes no truth. Stations with fewer
than four fully observed calendar years are excluded, because the
leave-one-year-out estimator protocol requires them. The same stations
contribute 2{,}422 station-years with at least 360 daily slots and at least one
missing day; those form the real-calendar population of Section 4.4 of the
article.

\section{Protocol history}

The study was pre-registered before the reserved networks were opened on
2026-09-23, and several rounds of repair followed. Three early repairs are
recorded here rather than silently applied.

\textbf{The incumbent's imputed limit was mis-stated.} An earlier version placed
the comparison at a thirteen-week (91-day) limit attributed to Johnson et al.
(2021). The source says twenty weeks. The analysis code now uses 140 days.

\textbf{The weighting between outage arms was not declared.} The case set mixes
an empirical arm, drawn from the outage geometry the same networks experience,
with a stratified arm built to span the decision space. Their base rates differ
by an order of magnitude. The pre-registered test was run on the pooled set, as
declared, and the decomposition is reported in the earlier development round.

\textbf{The natural scale of the peak has two values.} Computed over the full
eligible roster it is 1.06\,$^{\circ}$C; recomputed inside the development case
table it is 1.08\,$^{\circ}$C. The primary tolerance is 0.50\,$^{\circ}$C
throughout; the tolerance sweep shows the axis ordering is unchanged between
0.25 and 1.0 standard deviations.

The 2026-10-19 round repaired two label-isolation defects: outer folds now keep
their own calibration scores, and networks that share a station are kept in one
folding unit. The 2026-10-20 round repaired the object mismatch and the tie rule
that this supporting information documents below. The round's analyses are
registered under \texttt{major\_revision\_20261020/}: the aligned admission
(\texttt{p0\_01\_aligned\_policy\_admission}), the selector calibration
(\texttt{p0\_02\_selection\_calibration}), the real-calendar adoption comparison
(\texttt{p0\_03\_real\_calendar\_adoption}), the hot-year closure
(\texttt{p0\_04\_hot\_year\_closure}), the repair attribution
(\texttt{p0\_05\_repair\_attribution}) and the figure sources
(\texttt{p0\_06\_figures}). Each holds a \texttt{SUMMARY.md}, a
\texttt{protocol\_fingerprint.json} with script and input hashes and the git
head, and the machine-readable tables the article cites.

\section{The object mismatch, and what it changed}

The 2026-10-19 run declared the stored policy reconstruction as the admitted
object, but built the failure event from the certificate's absolute error. The
two are different reconstructions of the same station-year. On the external
panel the policy error exceeds the tolerance in 1{,}996 of 13{,}017 cases and the
certificate error in 2{,}078; 408 cases fail only under the policy and 490 only
under the certificate, so the two events disagree on 898 cases, or 6.90\% of the
panel. On the development panel they disagree on 830 of 40{,}704 cases
(2.04\%; 478 policy-only, 352 certificate-only).

The aligned protocol recomputes everything---threshold selection, evaluation,
intervals and contrasts---on \texttt{abserr\_policy}, after checking in-script
that $\texttt{abserr\_policy}=|\texttt{hottest30}_{\text{policy choice}} -
\texttt{truth\_hottest30}|$ on every finite row. The certificate bounds survive
only as ranking scores.

Table~\ref{tab:si-repair} reports the 2$\times$2 consequence of the two repairs.
On the external panel the fitted-score advantage over the peak-window rule is
2.7 points when the event is the certificate and the tie rule admits whole blocks
(2.8 with jitter), and 7.2 points when the event is the policy reconstruction and
ties are admitted whole.
Most of the change (about 4.5 points) comes from measuring the right object; the
tie rule contributes about 1.3 points. On the development panel the two repairs
move the same contrast by less than one point, which is why the defect was easy
to miss without an external comparison.

\begin{table}
\caption{Retention at the 5\% cap under the two failure events and the two tie
rules. Same scores, same folds, same selection rule; entries are the mean of
three fold seeds. ``Jitter'' reproduces the 2026-10-19 construction and applies
it to calibration and evaluation alike.}
\label{tab:si-repair}
\tabletext
\centering
\begin{tabular}{llrrrr}
\hline
panel & event / tie & fitted & peak-window & length & fitted $-$ peak (pp) \\
\hline
external & policy / whole block & 0.739 & 0.667 & 0.396 & $+7.2$ \\
external & policy / jitter & 0.741 & 0.656 & 0.445 & $+8.5$ \\
external & certificate / whole block & 0.728 & 0.701 & 0.406 & $+2.7$ \\
external & certificate / jitter & 0.728 & 0.699 & 0.445 & $+2.8$ \\
development & policy / whole block & 0.704 & 0.607 & 0.305 & $+9.7$ \\
development & policy / jitter & 0.704 & 0.602 & 0.324 & $+10.2$ \\
development & certificate / whole block & 0.706 & 0.616 & 0.305 & $+9.0$ \\
development & certificate / jitter & 0.708 & 0.622 & 0.342 & $+8.6$ \\
\hline
\end{tabular}
\end{table}

\section{Measurement contract and coverage}

Every result is reported against the daily-value statistic its providers actually
publish (Table~\ref{tab:si-contract}). Of the 40{,}704 development cases,
74.4\,\% come from the five providers that document a daily mean (84 networks),
15.5\,\% are project aggregates of sub-daily or instantaneous values (14
networks), 8.3\,\% come from ten networks whose provider never states whether the
published daily value is a mean, and 1.7\,\% are Finnish morning reference
readings. The external wave-2 panel is 94.4\,\% documented daily means. A
station-level coverage table with first and last year, complete years, statistic
class, quality class and licence route for every station is in
\texttt{major\_revision\_20261014/p0\_06\_coverage\_recompute/final\_coverage\_table.csv}.

The aligned rules were re-selected and re-evaluated inside the
documented-daily-mean subset (Table~\ref{tab:si-subset}). Thresholds were chosen
on the subset's calibration rows only and the stored per-fold calibration arrays
were restricted the same way; nothing was refitted. The fitted score's advantage
over the peak-window rule is $+8.0$, $+7.3$ and $+7.3$ points on the external
subset and $+8.8$, $+8.7$ and $+10.4$ on the development subset, with every
interval excluding zero, so the ranking reported in the article does not depend
on the project-aggregate or undocumented providers. Its comparison with the
certificate bound stays interval-crossing on development ($-2.3$ to $-4.2$
points) and is positive but mostly small on the external subset
($+0.1$ to $+1.9$).

\begin{table}
\caption{The aligned rules re-selected inside the documented-daily-mean subset
(30{,}289 development cases, 12{,}282 external cases; shares of the full panel
0.744 and 0.944). Retention is relative to the subset. Contrasts are paired,
group-clustered and reported as fitted minus the named rule in points.}
\label{tab:si-subset}
\tabletext
\centering
\begin{tabular}{llrrrr}
\hline
panel & rule & retention & failure rate & vs peak (pp) & vs bound (pp) \\
\hline
external & fitted risk & 0.749 & 0.031 & $+7.5$ & $+1.0$ \\
external & peak-window & 0.674 & 0.032 & --- & --- \\
external & certificate bound & 0.739 & 0.032 & --- & --- \\
external & gap-length & 0.411 & 0.018 & --- & --- \\
development & fitted risk & 0.733 & 0.030 & $+9.3$ & $-3.0$ \\
development & peak-window & 0.640 & 0.037 & --- & --- \\
development & certificate bound & 0.764 & 0.036 & --- & --- \\
development & gap-length & 0.347 & 0.019 & --- & --- \\
\hline
\end{tabular}
\end{table}

\begin{table}
\caption{Daily-value statistic by provider, as documented by the provider, with
the class assigned to each panel case.}
\label{tab:si-contract}
\small
\centering
\begin{tabular}{p{0.16\textwidth}p{0.20\textwidth}rp{0.27\textwidth}p{0.28\textwidth}}
\hline
provider & statistic documented & stations & class used in the analysis & evidence \\
\hline
USGS & daily mean & 427 & documented daily mean & NWIS daily-values service; parameter 00010 \\
GKD Bayern & daily mean & 55 & documented daily mean & provider daily series documentation \\
LUBW & daily mean & 22 & documented daily mean & provider daily series documentation \\
NVE & daily mean & 57 & documented daily mean & HydAPI daily values \\
FOEN & daily mean & 11 & documented daily mean & opendata.swiss daily series \\
CHMI & sub-daily (10-minute) & 73 & project aggregate & provider sub-daily series \\
Rijkswaterstaat & raw measurements & 5 & project aggregate & WaterWebServices raw readings \\
Hub'Eau & continuous & 9 & project aggregate & Hub'Eau chronique \\
ARSO & daily values, statistic unstated & 47 & undocumented & provider page only \\
SYKE & 08:00 near-shore reading & 3 & non-daily reference reading & provider documentation \\
\hline
\end{tabular}
\tablenotetext{a}{Counts are station-network rows from the coverage table; a
station appearing in two networks is counted once per network. The class
assignment is a deterministic parse of
\texttt{measurement\_provenance\_dictionary.csv}.}
\end{table}

\section{Parameters and their roles}

Every quantity that enters a result has one of four roles, and mixing them is a
common source of confusion (Table~\ref{tab:si-params}).

\begin{table}
\caption{Model and decision parameters by role. ``Fitted'' means estimated from
data at run time; ``fixed'' means a constant chosen before any result was seen;
``calibrated'' means chosen on groups disjoint from the cases it is applied to;
``declared'' means a study tolerance rather than a model quantity.}
\label{tab:si-params}
\small
\centering
\begin{tabular}{p{0.30\textwidth}p{0.14\textwidth}p{0.50\textwidth}}
\hline
quantity & role & value or source \\
\hline
ridge penalty of the donor and air estimators & fixed & $\alpha = 1.0$, seven air lags, whole-year fit \\
gradient-boosted risk score & fitted & quantile 0.9, depth 4, 350 iterations, learning rate 0.06, leaf 40; refit per outer fold \\
risk-score target & declared & absolute error of the admitted policy reconstruction \\
admission threshold & calibrated & largest candidate whose independence-group clustered bootstrap percentile at $1-0.05/n_{\mathrm{cand}}$ stays under 5\,\% \\
tie rule & declared & no jitter; a threshold is an observed score and $s\le\tau$ admits whole tie blocks \\
consequential-error tolerance $\tau$ & declared & $0.5\,^{\circ}$C \\
fold unit & fixed & connected component of networks that share a station id \\
folds & fixed & 5 outer folds, 3 fold seeds (101/202/303), 4 inner folds for the risk score \\
bootstrap draws & fixed & 1{,}000 for threshold selection, 2{,}000 for evaluation intervals, 10{,}000 for the stability check \\
\hline
\end{tabular}
\end{table}

Software and data are archived as follows. The analysis code, the frozen
configuration, the derived tables and the figure sources are in the accompanying
repository; the study inputs are packaged for confidential review
(\texttt{dist/wrr\_study\_inputs\_v1.tar.gz}) and a permanent archive with a DOI
is deposited on acceptance. Restricted material is listed rather than shipped:
LUBW raw values are read locally and not redistributed, and providers whose terms
name no licence are used under attribution only
(\texttt{metadata/data\_rights.csv}). Input-file hashes for every network are in
\texttt{current\_data\_manifest.csv}, and each analysis writes a protocol
fingerprint holding the script hash, the input hashes and the git head.

\section{Retention and hot-year composition in one place}

Table~\ref{tab:si-cap} collects the retention and realised risk of every rule on
both designed panels, at the same 5\% cap, for readers who want the decision
table without the figure. The reference arm is included for scale: on the
designed panels it admits everything and therefore fails the cap by a wide
margin, which is why it does not appear in the article's panel (a).

Table~\ref{tab:si-hot} gives the hot-year composition behind Figure~5c. Under
the pre-declared unique-station-year weighting the fitted score's total
deviation is $+0.9$ points ($[0.0,1.8]$) rather than $+0.7$, the ordering of the
rules is unchanged, and the gap-length rule's refusal of true hot years remains
the largest single effect in the table, but that is its retention and not a
selection against hot years: against the retention of all records the excess
refusal of true hot years is negative for every rule on the external panel
(Table~\ref{tab:si-hotnull}). The case-weighted estimate is primary because it is
the archive an agency holds under this monitoring design.

\section{The selector's risk control, read correctly}

The deployed threshold reads the percentile $1-0.05/n_{\mathrm{cand}}$ of an
independence-group clustered bootstrap of the calibration false-admit rate, with
$n_{\mathrm{cand}}$ the number of distinct snapped thresholds. Two properties
were separated in simulation on 2{,}000 synthetic clustered archives per
scenario, with the deployed rule imported from the analysis code so the procedure
cannot drift:
\begin{itemize}
\item the \emph{population} risk of the selected threshold, measured on an
independent two-million-case reference draw from the same generator; and
\item the error rate observed in a finite fresh test batch of 16 groups.
\end{itemize}
Table~\ref{tab:si-selector} reports both. With the deployed simultaneous bound,
the selected threshold's population risk exceeds the 5\% cap in 0.4\% of
repetitions in the baseline scenario, 0.6\% near the cap, 0.4\% with unequal
group sizes and 2.4\% with twelve calibration groups. The finite test batch
exceeds the cap far more often (7.8\%, 7.5\%, 4.6\% and 12.2\%) because the batch
adds noise the rule cannot control. The 8.5\% quoted in the previous version of
this work was the second quantity; reading it as the selector's error rate was a
misinterpretation. A pointwise bound keeps about twice as many records (35.9\%
against 18.5\% in the baseline scenario) and violates the cap about six times as
often (7.4\% against 0.4\% on the population measure).

\begin{table}
\caption{Selector calibration per scenario, 2{,}000 repetitions each. ``True
risk'' is the population failure rate of the selected threshold on an independent
two-million-case reference draw; ``batch'' is the rate observed in a finite
fresh test sample.}
\label{tab:si-selector}
\tabletext
\centering
\begin{tabular}{llrrrrr}
\hline
scenario & procedure & retention & true risk & P(true $>5\%$) & P(batch $>5\%$) & batch rate \\
\hline
baseline & deployed & 0.185 & 0.0199 & 0.004 & 0.078 & 0.0204 \\
baseline & pointwise & 0.359 & 0.0352 & 0.074 & 0.191 & 0.0356 \\
low error & deployed & 0.761 & 0.0313 & 0.000 & 0.051 & 0.0320 \\
low error & pointwise & 0.881 & 0.0386 & 0.000 & 0.153 & 0.0394 \\
near cap & deployed & 0.288 & 0.0222 & 0.006 & 0.075 & 0.0224 \\
near cap & pointwise & 0.543 & 0.0365 & 0.066 & 0.200 & 0.0369 \\
unequal groups & deployed & 0.296 & 0.0306 & 0.004 & 0.046 & 0.0308 \\
unequal groups & pointwise & 0.425 & 0.0406 & 0.064 & 0.181 & 0.0408 \\
few groups & deployed & 0.153 & 0.0171 & 0.024 & 0.122 & 0.0173 \\
few groups & pointwise & 0.309 & 0.0311 & 0.119 & 0.254 & 0.0322 \\
\hline
\end{tabular}
\end{table}

The extreme percentile is numerically delicate at 1{,}000 draws. Re-running the
selection with 10{,}000 draws changes the chosen threshold in 42 of 150
fold-by-score cells. The fitted score's retention moves by at most 1.0 point and
the discrete rules by at most 3.9 points. The article therefore describes the
selected threshold as a calibration decision at this granularity, not as a stable
optimum.

\section{Predictions that failed, and one that was withdrawn}

\textbf{The prediction that failed this round.} Before the aligned rerun we
predicted, on the basis of the 2026-10-19 numbers, that the fitted score would
not clear the pre-specified five-point margin against the peak-window rule even
after the object was repaired. It clears it on the external panel at all three
fold seeds (point estimates 8.0, 7.3 and 7.0 points, lower interval limits 5.0, 4.6
and 4.3) and on the development panel (9.7 to 10.3 points). The earlier reading
was an artefact of the object mismatch, and we report the corrected result rather
than the prediction.

\textbf{A justification that did not survive the ladder.} The article's earlier
account of the fitted score's gain was that it reads something the support share
does not contain. That is true, but the 2026-10-14 same-information ablation had
already found that adding the support share to a generic model changes retention
by $0.00$ points ($[-1.00, +1.00]$), and the manuscript did not carry it. The
aligned protocol reproduces it (Table 3 of the article and Tables~\ref{tab:si-ladder-ext}
and \ref{tab:si-ladder-dev} here): the support share adds $+0.3$ to $+0.8$ points on the
external panel and $-0.5$ to $+0.5$ on the development panel, and the gain over the
peak-window rule belongs to donor information and reconstruction diagnostics.

\textbf{A claim withdrawn rather than falsified.} An earlier version framed the
result as ``position, not length, governs''. The crossed design does not support
a law: the support effect is universal in sign but its size ranges from 0.24 to
0.64\,$^{\circ}$C at 30 days across reconstructors, and linear interpolation
retains a large length effect at a high hidden share (0.82\,$^{\circ}$C
[0.64, 1.00] between 90 and 30 days at a true share of at least 0.9). The other
four reconstructors keep a small length effect, 0.05 to 0.14\,$^{\circ}$C, that
is distinguishable from zero (Table~\ref{tab:si-length-estimators}). The article now states the conditional version.

\textbf{Three earlier predictions that were not met} remain in the record from
the development round: the support-restricted advantage does not decay
monotonically in support width; the estimand's own support is not the best
conditioning set at every width; and two hand-chosen phase windows do not beat
both peak overlap and outage length.

\section{Runs that were discarded}

Five defects made specific artifacts unusable, and each was repaired before the
numbers reported here:
\begin{enumerate}
\item outage realisations were once seeded from a process-salted hash, so the case
  table could not be regenerated;
\item worker results were once concatenated in completion order, so row order was
  not reproducible;
\item an additive per-estimator conformal scheme under-covered, because selecting
  the smallest of several bounds is a winner's curse that per-estimator
  calibration cannot see;
\item a shared calibration-score array let a later outer fold overwrite earlier
  folds' scores with models that had seen those folds' labels
  (repaired 2026-10-19);
\item the failure event was read from the certificate reconstruction while the
  admitted object was the policy reconstruction, and the two tie rules differed
  between calibration and deployment (repaired 2026-10-20).
\end{enumerate}
In each case the affected numbers were discarded rather than reported alongside
the repaired ones.

\section{The attribution round: design audit, ladder, channels and rule}

This section holds the tables behind the article's attribution claims. They are
written by \texttt{scripts/346\_wrr\_si\_tables\_r2.py} from the artifacts of
\texttt{scripts/340} to \texttt{scripts/344} under
\texttt{results/wrr\_submission\_ready/major\_revision\_20261020b/}, and nothing in
them is typed by hand.

\textbf{What the designs contain.} The earlier article described the 53{,}721
designed outages as a factorial design with scattered extra missing days. The case
tables say otherwise, and Tables~\ref{tab:si-arms} to \ref{tab:si-cells} give the
corrected description: single contiguous runs of 14, 30, 60, 90, 120, 180 or 270
days (stratified) and transplanted real geometries (empirical); no scattered days;
the overlap class defined on the climatological window; and a correlation of gap
length with the true hidden share of $+0.52$ that the design does not remove. The
crossed panel (1{,}899 outages) was generated separately on 24 development stations;
it also holds 19 transplanted outages, which no analysis here uses.

\textbf{The ladder.} Tables~\ref{tab:si-ladder-levels}, \ref{tab:si-ladder-ext} and
\ref{tab:si-ladder-dev} give the retention and the paired contrasts of every
specification of the information ladder, including the two non-deployable
specifications that read the true hidden share.

\textbf{Channels, interventions and the misalignment contrast.}
Tables~\ref{tab:si-channels} to \ref{tab:si-misalign-class}.

\textbf{Sensitivity of the length effect.} Table~\ref{tab:si-length-estimators}
compares the within-station-year estimator of the article with the network-subset
estimator of the earlier mechanism run; Table~\ref{tab:si-competing} gives the
competing-peak relation.

\textbf{The screening-need rule and the hot-year null.} Tables~\ref{tab:si-need}
to \ref{tab:si-hotcost}. The real replay's evidence is two to four networks that
need a screen on 3 to 6 target station-years each, and the tables are the
complete record of it.

\input{si_r2_tables}

\section{Verification and reproduction}

\begin{itemize}
  \item The case tables are byte-reproducible including row order, pinned by a
    test.
  \item No feature may read a day the outage hides; enforced in the feature
    constructor and pinned by a test.
  \item The aligned admission script asserts that \texttt{abserr\_policy} equals
    the absolute difference between the stored policy prediction and the
    reference value before it uses it, so an alignment error aborts the run.
  \item The failure event, the threshold and the evaluated error are the same
    column in every stage of the aligned protocol; a unit test checks that the
    case-level admissions reproduce the thresholds.
  \item Every table cell and audited inline number in the article is re-derived
    from the artifacts by the build script, which fails when a listed number is
    missing from the text.
  \item Every input the current study reads is listed with a byte count and a
    SHA-256 in \texttt{current\_data\_manifest.csv}, and each analysis records
    its own protocol fingerprint.
\end{itemize}

\begin{table}
\caption{Retention at the 5\% cap by rule and panel (mean of three fold seeds,
with the seed range in parentheses), and the realised failure rate among admitted
records. The reference arm reconstructs every record and is included for scale;
it does not satisfy the cap on the designed panels.}
\label{tab:si-cap}
\tabletext
\centering
\begin{tabular}{lrrrr}
\hline
rule & external retention & external risk & development retention & development risk \\
\hline
gap-length & 0.396 & 0.018 & 0.305 & 0.024 \\
peak-window & 0.667 & 0.033 & 0.604 & 0.038 \\
certificate bound & 0.729 & 0.032 & 0.727 & 0.036 \\
fitted risk & 0.740 & 0.031 & 0.704 & 0.030 \\
reconstruct everything & 1.000 & 0.153 & 1.000 & 0.149 \\
\hline
\end{tabular}
\end{table}

\begin{table}
\caption{Hot-year composition, external panel, 20\,$^{\circ}$C, fold seed 101,
case-weighted. Shares are fractions. The unique-station-year weighting is
reported in the text and leaves the ordering unchanged.
``Refused hot'' is the share of true hot years left out of the archive;
``missed hot'' is the share the admitted reconstruction places below the
threshold.}
\label{tab:si-hot}
\tabletext
\centering
\begin{tabular}{lrrrrrr}
\hline
rule & p$_{\text{full}}$ & p$_A$ & selection & total (95\%) & refused hot & missed hot \\
\hline
peak-window & 0.857 & 0.861 & $+0.004$ & $+0.003$ $[-0.005,0.012]$ & 0.332 & 0.003 \\
fitted risk & 0.857 & 0.868 & $+0.011$ & $+0.007$ $[-0.002,0.016]$ & 0.245 & 0.004 \\
certificate bound & 0.857 & 0.868 & $+0.011$ & $+0.007$ $[-0.002,0.017]$ & 0.262 & 0.004 \\
gap-length & 0.857 & 0.875 & $+0.018$ & $+0.016$ $[-0.003,0.046]$ & 0.595 & 0.001 \\
reconstruct everything & 0.857 & 0.857 & $0.000$ & $-0.006$ $[-0.013,0.003]$ & 0.000 & 0.013 \\
\hline
\end{tabular}
\end{table}

\begin{figure}
\includegraphics[width=\textwidth]{figures/s1_study_design.pdf}
\caption{The evidence chain. Complete station-years are the reference truth; the
development and external panels place stratified (single-run) and transplanted
(real-geometry) outages and supply the admission comparisons; the real archive
supplies gap geometry only, because it has no complete counterpart; and the 418
real-calendar groups, from 59 recorded gaps on 67 target station-years, supply the
deployment probe. The separately generated crossed panel (1{,}899 outages on 211
station-years) that identifies the length and position effects is not drawn. Numbers give cases (and, for
the panels, independence groups). The external panel is already opened, so its
comparisons are cross-fitted within it rather than a sealed confirmation.
Schematic source:
\texttt{s1\_study\_design.drawio} (draw.io; see the Open Research statement of
the article).}
\label{fig:si-design}
\end{figure}

\end{document}
